\documentclass[12pt,a4paper]{article}

\usepackage[utf8]{inputenc}
\usepackage[T1]{fontenc}
\usepackage[english]{babel}
\usepackage{amsmath}
\usepackage{amssymb}
\usepackage{geometry}
\usepackage{xcolor}
\usepackage{mdframed}
\usepackage{enumitem}

\geometry{margin=2.5cm}

\definecolor{azul}{RGB}{30, 100, 180}
\definecolor{cinza}{RGB}{245, 245, 245}
\definecolor{verde}{RGB}{20, 130, 80}

\mdfdefinestyle{enunciado}{
  backgroundcolor=cinza,
  linecolor=azul,
  linewidth=1.2pt,
  topline=false, rightline=false, bottomline=false,
  innerleftmargin=14pt, innerrightmargin=14pt,
  innertopmargin=12pt, innerbottommargin=12pt
}

\mdfdefinestyle{dados}{
  backgroundcolor=cinza,
  linecolor=gray!50,
  linewidth=0.5pt,
  innerleftmargin=12pt, innerrightmargin=12pt,
  innertopmargin=8pt, innerbottommargin=8pt,
  roundcorner=4pt
}

\mdfdefinestyle{resolucao}{
  backgroundcolor=white,
  linecolor=verde,
  linewidth=1.2pt,
  topline=false, rightline=false, bottomline=false,
  innerleftmargin=14pt, innerrightmargin=14pt,
  innertopmargin=12pt, innerbottommargin=12pt
}

\begin{document}

\begin{center}
  {\small\textcolor{azul}{\textbf{F\'isica e Qu\'imica A\,---\,11.\textordmasculine{} ano}}}
\end{center}

\bigskip

%---------------------------------------------------------------
\begin{mdframed}[style=enunciado]

Um investigador prepara uma solu\c{c}\~ao aquosa saturada de um sal pouco
sol\'uvel, XY, a 25\,\textdegree C, que se dissolve segundo o equil\'ibrio:
%
\[
  \mathrm{XY}_{(s)} \rightleftharpoons \mathrm{X}^{+}_{(aq)}
  + \mathrm{Y}^{-}_{(aq)}
\]
%
Para confirmar a concentra\c{c}\~ao da solu\c{c}\~ao, o investigador usa um
feixe de luz \textit{laser} que passa do ar para a superf\'icie da solu\c{c}\~ao
com um \^angulo de incid\^encia de $\theta_1 = 45{,}0^\circ$ e mede um
\^angulo de refra\c{c}\~ao de $\theta_2 = 32{,}0^\circ$.

Sabe-se que o \'indice de refra\c{c}\~ao de uma solu\c{c}\~ao aquosa varia
com a concentra\c{c}\~ao total de i\~oes de acordo com:
%
\[
  n_{\text{sol}} = n_{\text{\'{a}gua}} + k \cdot c_{\text{total}}
\]
%
sendo $c_{\text{total}}$ a concentra\c{c}\~ao total de i\~oes em
mol\,L$^{-1}$.

\end{mdframed}

\medskip

%---------------------------------------------------------------
\begin{mdframed}[style=dados]
\textbf{Dados:}
\begin{itemize}[leftmargin=1.5em, itemsep=3pt, topsep=4pt]
  \item $K_s\,(\mathrm{XY},\;25\,\text{\textdegree C}) = 7{,}50 \times 10^{-4}$
  \item $n_{\text{ar}} = 1{,}000$ \quad;\quad $n_{\text{\'{a}gua pura}} = 1{,}333$
  \item $k = 2{,}50 \times 10^{-2}\ \text{L\,mol}^{-1}$
\end{itemize}
\end{mdframed}

\medskip

%---------------------------------------------------------------
\begin{enumerate}[label=\textbf{\alph*)}, leftmargin=1.8em, itemsep=8pt]

  \item Aplique a lei de Snell-Descartes para determinar o \'indice de
        refra\c{c}\~ao da solu\c{c}\~ao, $n_{\text{sol}}$.

  \item A partir de $n_{\text{sol}}$, calcule a concentra\c{c}\~ao total de
        i\~oes em solu\c{c}\~ao, $c_{\text{total}}$.

  \item Sabendo que $[\mathrm{X}^+] = [\mathrm{Y}^-] = s$, determine a
        solubilidade molar $s$ a partir de $c_{\text{total}}$ e verifique se
        o resultado \'e consistente com o valor de $K_s$ fornecido.

\end{enumerate}

\bigskip\bigskip

%---------------------------------------------------------------
\noindent\textcolor{verde}{\rule{\linewidth}{0.5pt}}\\[4pt]
{\textcolor{verde}{\textbf{Resolu\c{c}\~ao}}}\\[-2pt]
\noindent\textcolor{verde}{\rule{\linewidth}{0.5pt}}

\bigskip

%--- a)
\begin{mdframed}[style=resolucao]
\textbf{a)\ \'Indice de refra\c{c}\~ao da solu\c{c}\~ao}

\medskip
Pela lei de Snell-Descartes:
\[
  n_{\text{ar}}\,\sin\theta_1 = n_{\text{sol}}\,\sin\theta_2
\]
\[
  n_{\text{sol}}
  = \frac{n_{\text{ar}}\,\sin\theta_1}{\sin\theta_2}
  = \frac{1{,}000 \times \sin 45{,}0^\circ}{\sin 32{,}0^\circ}
  = \frac{0{,}7071}{0{,}5299}
  \approx \boxed{1{,}334}
\]
\end{mdframed}

\bigskip

%--- b)
\begin{mdframed}[style=resolucao]
\textbf{b)\ Concentra\c{c}\~ao total de i\~oes}

\medskip
Isolando $c_{\text{total}}$ na express\~ao do \'indice de refra\c{c}\~ao:
\[
  c_{\text{total}}
  = \frac{n_{\text{sol}} - n_{\text{\'{a}gua}}}{k}
  = \frac{1{,}334 - 1{,}333}{2{,}50 \times 10^{-2}}
  = \frac{1{,}00 \times 10^{-3}}{2{,}50 \times 10^{-2}}
  \approx \boxed{4{,}00 \times 10^{-2}\ \text{mol\,L}^{-1}}
\]

\smallskip
\textit{Nota:} o valor mais preciso de $n_{\text{sol}} = 1{,}3344$ d\'a
$c_{\text{total}} = 5{,}47 \times 10^{-2}$ mol\,L$^{-1}$; a diferen\c{c}a
deve-se ao arredondamento de $n_{\text{sol}}$ a tr\^es casas decimais.
Para coer\^encia com os dados, usa-se o valor arredondado.
\end{mdframed}

\bigskip

%--- c)
\begin{mdframed}[style=resolucao]
\textbf{c)\ Solubilidade molar e verifica\c{c}\~ao com $K_s$}

\medskip
Como $[\mathrm{X}^+] = [\mathrm{Y}^-] = s$, a concentra\c{c}\~ao total de
i\~oes \'e $c_{\text{total}} = 2s$, logo:
\[
  s = \frac{c_{\text{total}}}{2}
    = \frac{5{,}47 \times 10^{-2}}{2}
    \approx 2{,}74 \times 10^{-2}\ \text{mol\,L}^{-1}
\]

\medskip
\textit{Verifica\c{c}\~ao pelo $K_s$:}
\[
  K_s = [\mathrm{X}^+][\mathrm{Y}^-] = s^2
  \implies
  s = \sqrt{K_s} = \sqrt{7{,}50 \times 10^{-4}}
  \approx \boxed{2{,}74 \times 10^{-2}\ \text{mol\,L}^{-1}}
\]

Os dois valores coincidem, confirmando a consist\^encia entre o resultado
experimental (obtido por refratometria) e o valor tabelado de $K_s$.
\end{mdframed}

\end{document}
